
If the primal is a minimization problem, and the dual is written as a
maximization problem, then the following rules apply. 
\begin{itemize}
\item For each constraint of the primal there is a dual variable. The
  sign of the dual variable depends on the sign of the primal constraint, according to the following table:
\begin{center}
\begin{tabular}{c | c}
Constraint of the primal & Dual variable \\
\hline
$=$ & free\\
$\leq$ & $\leq 0$\\
$\geq$ & $\geq 0$
\end{tabular}
\end{center}
\item For each primal variable there is a dual constraint. The sign of
  the constraint depends on the sign of the primal variable, according to the table below.
\begin{center}

\begin{tabular}{c | c}
Primal variable & Dual constraint \\
\hline
$\geq 0$ & $\leq$\\
$\leq 0$ & $\geq$\\
free & $=$
\end{tabular}
\end{center}

\end{itemize}


Write the dual of the following linear optimization problem in two
ways: first use the summary tables above, and then use the Lagrangian
method in order to see the correctness of those summary tables.
\[
\min   x_1 - 4x_2 + 2x_3 
\]
subject to
\begin{align*}
 2x_1 + x_2 &\geq 10, \\
x_1 + 4x_3 &= 16, \\
x_2 - x_3 &\leq 5, \\
x_1 &\geq 0, \\
x_2 &\in \mathbb{R}, \\
x_3 &\leq 0.
\end{align*}

